\section{A descending-binomial quartic in 233 natural variables}
\label{sec:natural}

This branch starts from the direct numerical-calculus construction of
package 129 \cite{parent129}. It changes the factorial certificate and
the associated reciprocal CRT moduli, combines the semantic anchor with
a single cap equation, and uses a checked horizon floor of 12. It
retains the Cantor-coded calculus, six exact power requests, eleven
modular binomial requests, and the dominant denominator. Its final
polynomial has 233 variables, 182 residuals, 7,313 collected monomials,
and height 86,543,666.

\subsection{A short factorial fixed point}

The rising-binomial parent used a ceiling and had to exclude spurious
small values before obtaining a strong lower bound for the factorial
witness. Reversing the product gives a particularly direct lower bound.

\begin{lemma}[Descending factorial identity]\label{lem:floor}
Let \(t\ge3\), \(F\ge2\), and \(F^2-1\ge t\) be integers. Then
\[
 F=\left\lfloor\frac{(F^2-1)^t}{\binom{F^2-1}{t}}\right\rfloor
 \quad\Longleftrightarrow\quad F=t!.
\]
\end{lemma}

\begin{proof}
Put \(n=F^2-1\), \(T=t!\), and \(S=t(t-1)/2\). The ratio is
\[
 H=T\prod_{j=0}^{t-1}\left(1-\frac{j}{n}\right)^{-1}\ge T.
\]
If its floor is \(F\), then \(F\ge T\). For \(t\ge3\),
\(T-1>S\); this is immediate at \(t=3\) and persists by induction.
Consequently
\[
 n\ge T^2-1=(T+1)(T-1)>(T+1)S.
\]
The elementary product inequality
\(\prod_j(1-a_j)\ge1-\sum_j a_j\), for \(0\le a_j\le1\), gives
\[
 T\le H\le\frac{Tn}{n-S}
       =T+\frac{TS}{n-S}<T+1.
\]
Thus the floor is \(T\), so \(F=T\). Conversely, if \(F=T\),
the same displayed inequalities apply and the floor is \(F\).
\end{proof}

The actual polynomial graph makes \(F=f_0+2\), with \(f_0\in\N\),
and evaluates
\[
 n=(f_0+1)(f_0+3)=F^2-1.
\]
This expression has nonnegative coefficients. Its two floor rows are
\begin{equation}\label{eq:floor-rows}
 V_0=BF+r,\qquad B=r+s+1,\qquad r,s\in\N.
\end{equation}
After the arithmetic interfaces have recovered
\(V_0=n^t\) and \(B=\binom nt\), these rows say exactly
\(BF\le V_0<B(F+1)\). The literal shift \(B=b_0+2\) ensures
\(B\ge2\), so the recovered binomial is nonzero and \(n\ge t\).
Lemma~\ref{lem:floor} then gives \(F=t!\). No factorial, floor, or
binomial primitive survives in the final polynomial.

\subsection{Descending reciprocal moduli}

\begin{lemma}\label{lem:descending-crt}
Let \(F=t!\) with \(t\ge3\). For \(1\le i\le t\), put
\[
 d_i=\frac{F^2}{i}-1.
\]
These are positive integers greater than one, are pairwise coprime,
have no prime divisor at most \(t\), and satisfy
\[
 \prod_{i=1}^{t}d_i=\binom{F^2-1}{t}=B.
\]
\end{lemma}

\begin{proof}
Each \(i\mid F\), hence \(F^2/i\) is a multiple of \(F\) and
\(d_i\equiv-1\pmod F\). Every prime at most \(t\) divides \(F\),
so no such prime divides \(d_i\). If a prime divided both \(d_i\)
and \(d_j\), it would divide
\(i d_i-jd_j=j-i\), which is impossible for distinct indices in
\(\{1,\ldots,t\}\). The product identity is
\[
 \prod_{i=1}^{t}\frac{F^2-i}{i}
 =\frac{(F^2-1)\cdots(F^2-t)}{t!}.
\]
\end{proof}

The encoded local-row index \(C\) now uses the polynomial equality
\begin{equation}\label{eq:descending-index}
 C+1=F^2+Bq_C,\qquad q_C\in\N,
\end{equation}
in place of the parent's \(C+F^2+1=Bq_C\). Modulo \(d_i\),
\(F^2\equiv i\), so the recovered local index is still \(i-1\).
The quotient's natural domain is legitimate in completeness. Indeed,
\(B>F^2\): with \(n=F^2-1\), we have
\(\binom nt\ge\binom n2>n+1\) in the present range. Every positive
CRT representative of \(F^2-1\pmod B\) is therefore at least
\(F^2-1\), and \(q_C\ge0\).

An additional consequence is \(d_i\equiv-1\pmod4\), because
\(4\mid F\) in the actual \(t\ge12\) construction. Thus each
\(d_i\) has a prime divisor congruent to 3 modulo 4 with odd
multiplicity. This observation is not used to claim another compression
in the delivered equation; the preceding elementary CRT facts suffice.

\subsection{One cap equation at a semantic index}

For \(A\ge2\), write
\[
 X_j(A)+Y_j(A)\sqrt{A^2-1}=(A+\sqrt{A^2-1})^j.
\]
The nonnegative solutions of
\(x^2-(A^2-1)y^2=1\) are exactly these pairs, with unique \(j\ge0\).
The elementary recurrences give
\[
 Y_j(A)\equiv j\pmod{A-1},\quad Y_j(A)\ge j,\quad
 Y_{j+1}(A)>(2A-1)Y_j(A)\quad(j\ge1).
\]
The complete one-index certificate of the parent is retained for a
single pair \((X_*,Y_*)=(X_H(A),Y_H(A))\). The new simultaneous
constraint is
\begin{equation}\label{eq:natural-one-row}
 \boxed{\quad
 Y_*=H+\sum_{i=1}^{r}(y_i-k_i)^2+(A-1)Q,
 \qquad Q\in\N.\quad}
\end{equation}
It combines the size cap and the error congruence in one quadratic row.
Before the anchor is interpreted, this row already gives \(H\le Y_*\).
It therefore supplies the index bound required by the retained
one-index certificate.
This is a specialization and combination of the earlier one-row cap
idea \cite{cap128} with package 129's semantic anchor, not a priority
claim for the general idea of combining a cap and a congruence.

\begin{lemma}\label{lem:natural-cap}
Assume \(W\ge3\), \(2\le H<W\), \(r\le W-2\),
\(0\le k_i<W\), \(A>W+1\), and \(A-1>W^3\).
Let \(Y_*=Y_H(A)\), and let each \((x_i,y_i)\) be a
nonnegative Pell solution at base \(A\). If
\eqref{eq:natural-one-row} holds, then
\((x_i,y_i)=(X_{k_i}(A),Y_{k_i}(A))\) for every \(i\).
Conversely, for the positive indices in this construction, the
canonical pairs satisfy \eqref{eq:natural-one-row} with a natural
\(Q\) whenever \(H>2\max_i k_i\).
\end{lemma}

\begin{proof}
Write \(y_i=Y_{j_i}(A)\). If \(j_i\ge H\), then
\(y_i\ge Y_*\ge Y_2(A)=2A>2W\), so
\[
 (y_i-k_i)^2\ge(Y_*-W+1)^2>Y_*.
\]
This contradicts the nonnegative right side of
\eqref{eq:natural-one-row}. Therefore \(j_i<H<W\).
Reducing the row modulo \(A-1\) gives divisibility of
\[
 E=\sum_i(j_i-k_i)^2,
 \qquad 0\le E\le r(W-1)^2<W^3<A-1.
\]
Hence \(E=0\), proving every index equality.

For completeness put \(m=\max_i k_i\ge1\). The identity
\[
 Y_{2m+1}(A)=Y_{m+1}(A)^2-Y_m(A)^2
            >4A(A-1)Y_m(A)^2
\]
and \(H\ge2m+1\) show
\[
 Y_H(A)-H-\sum_i(Y_{k_i}(A)-k_i)^2
 >[4A(A-1)-r]Y_m(A)^2-H>0.
\]
This integer is divisible by \(A-1\) by the recurrence congruence.
Its quotient is the required nonnegative \(Q\).
\end{proof}

The anchor is the existing derived label polynomial \(H=J\), which
includes the Cantor constructor
\[
 I(a,b)=2(a+b)(a+b+1)+4b,
 \qquad J(a,b,c,d)\ge I(a,b)\ge2a^2.
\]
As in the parent, the nine independent CRT representatives initially
lie in \([1,N^2]\). Increase the encoded \(a\)-column by \(N^2\).
All local polynomial congruences at precision \(d_i^{2K}\) are
preserved, and their aggregate remains at most \(10N^2\).
The new \(a>N^2\) gives \(H>2N^4\). Every other required index
is at most \(2N^2\), so \(H>2\max_i k_i\), as required for the
new converse. No extra dominance inequality appears in the polynomial.
Soundness uses only the cap lemma's explicit frame hypotheses.

\subsection{The six exact powers and the order of recovery}

The semantic parameters are
\[
 t=t_0+12,\quad F=f_0+2,\quad
 R=2+t+x+y+F,\quad K=2R,\quad B=b_0+2.
\]
The six exact requests are
\begin{equation}\label{eq:natural-six}
 V_0=(F^2-1)^t,\qquad N=B^K,\qquad q=N^K,\qquad
 D_t=q^t,\quad D_y=q^{y+1},\quad D_x=q^{x+1}.
\end{equation}
These are meanings proved for finite polynomial graphs, not operators
allowed in \(P_{\N}\). The extra denominator \(q^{t+1}=qD_t\)
uses a multiplication already accounted for in the graph.

The unchanged enlarged growth frame is
\[
 W=16+\sum_{a\in\mathcal A}a,\quad Z=z+1,\quad
 L^2-\big((W+1)^2-1\big)(WZ)^2=1,\quad A=L+\sigma,
\]
where \(\mathcal A\) contains every distinct base, exponent, and
proposed output in \eqref{eq:natural-six}, as well as every modular
upper index. All these atoms are natural and less than \(W\).
Pell classification at base \(W+1\), followed by reduction modulo
\(W\), shows that its nonzero index is at least \(W\). Therefore
\[
 A\ge L\ge X_W(W+1)\ge(W+1)^W,
 \qquad A-1>W^3.
\]
The actual base \(A\) remains free to increase after the lower witness
\(L\) is chosen. This freedom is needed for the later root congruence.

There are fourteen ordinary index pairs and one fully certified
anchor. For ordinary pairs the graph uses consecutive ordinates
\(y,z\) and a product wire \(u\):
\[
 u=yz,\qquad y^2+z^2-2Au=1.
\]
Every such pair has an output row of the form
\(ny=z+R_0\), with \(R_0\ge0\) and \(1\le n<A\).
Thus \(y>0\), \(z<Ay\), and \(x=Ay-z>0\) gives an ordinary
Pell pair. The one-row cap authenticates all its indices before
interpreting a power or binomial coefficient.

For a proposed power \(m=n^k\), put
\(T_n=2An-n^2-1\) and impose
\[
 X_k(A)+(n-A)Y_k(A)=m+T_n h_n,\qquad h_n\in\N.
\]
Modulo \(T_n\), the left side is \(n^k\). The frame proves
\(T_n>W^W+W\), \(m<W\), and \(n^k<W^W\), so both small
residues must agree. This proves all six exact powers. The canonical
left side is positive and \(0\le m<T_n\), so their difference cannot
be a negative multiple of \(T_n\). Its congruence quotient is therefore
natural in completeness. The literal modulus row also excludes a zero
power base.

Before using any factorial semantics, \(B\ge2\) gives
\[
 q=B^{K^2}\ge2^{4R^2}.
\]
With \(n=F^2-1<R^2\), every coefficient \(\binom nj\),
\(0\le j\le t\), is less than
\(R^{2t}\le2^{2Rt}<2^{2R^2}<q\). Thus the first binomial digit
can be decoded before \(F=t!\) is known. Its output is \(B\);
the floor rows now recover \(F=t!\) by Lemma~\ref{lem:floor}.
There is no circular assumption that a large proposed factorial is
already the true factorial.

\subsection{The common root and exact digits}

The unchanged endpoint rows include
\[
 y=F+1+r_0,\quad
 x=F+1+(y+1)u_0,\quad
 x=1+\bigl(1+(t+1)y\bigr)u_t.
\]
They imply \(x\ge2+(t+1)y>\max(t+1,y)\). Consequently
\(M=ND_x=Nq^{x+1}\) is a multiple of every needed digit modulus.
Put \(b=q+1\) and choose the common-root equation
\[
 2Ab=b^2+1+M h_0.
\]
It implies
\[
 \rho_n=X_n(A)+(b-A)Y_n(A)\equiv b^n\pmod M.
\]
For an exact binomial output \(C=\binom nk\), with \(d=q^k\),
the literal rows are
\[
 \rho_n=r+dC+qd h,\qquad d=r+s+1,\qquad q=C+o+1.
\]
For a range request \(N\mid\binom nk\), they are
\[
 \rho_n=r+Nd h,\qquad d=r+s+1.
\]
The carry-free binomial theorem gives precisely these interpretations.
For completeness, \(\rho_n>0\), while \(0\le r+dC<qd\) for an
exact digit and \(0\le r<d\le Nd\) for a range digit. Subtracting
the indicated small residue therefore cannot give a negative multiple
of the corresponding modulus, so the reconstructed high-digit
quotients are nonnegative. Also \(N\mid q\), so \(\gcd(b,M)=1\).
The linear congruence \(2bA\equiv b^2+1\pmod M\) is soluble:
its only possible common divisor is 2, and then \(b^2+1\) is even.
There are arbitrarily large suitable bases \(A\ge L\).

The remaining ten range binomials recover the eight local source
columns, the row index, and the derived label sum. Their upper-index
and carry bounds are the parent's bounds with the descending factorial
input, which is smaller than the former \(F^2+t\). In completeness,
the semantic-anchor lift still uses the same \(10N^2\) aggregate.
For representatives \(n\le10N^2\), relevant lower indices
\(k\le R-1\), and \(0\le j\le k\), the old estimate gives
\(\binom nj<q\).
For the derived label sum, the exact source certificate gives
\(J\le(n+1)^{\ell}\), \(\ell=2^{22}\), and
\(3\ell(t+1)\le2F\le K\) gives the required stronger bound.

\subsection{Why horizon 12 is sufficient}

The literal majorant certificate has local exponent \(E=2^{23}\)
and label exponent \(\ell=2^{22}\). At the new floor,
\[
 12!=479001600,\qquad a=\lfloor\log_2(12!)\rfloor=28,
\]
and direct integer arithmetic verifies
\[
 aE<2^a,\qquad 2E\le2^a,\qquad
 3\ell\cdot13\le2\cdot12!.
\]
These imply the same unbounded inequalities used by the source.
In particular, \(R^E<2^R\) starts at \(R=2^a\); the ratio of
successive left sides is below 2 for \(R\ge2E\), so it persists
for larger integers. The label inequality persists because a factorial
grows faster than its linear comparison. Thus the smaller literal
horizon changes coefficients without weakening a soundness estimate.

\subsection{From the local arithmetic rows to refutations}

For clarity, the inherited CRT and logical steps are stated explicitly.
If \(p^e\mid\binom ak\), then \(a\bmod p^e<k\). For
\(a\ge k\), this follows from the valuation formula for binomials:
if the residue were at least \(k\), the valuation contributions at
levels \(e,e+1,\ldots\) would vanish, leaving valuation at most
\(e-1\). The case \(a<k\) is immediate. The converse is not
assumed.

For each \(d_i\), choose a prime divisor \(p\), let
\(m=v_p(d_i)\), and work modulo \(p^{Km}\). The index congruence
\eqref{eq:descending-index} and its range binomial recover \(i-1\).
The other range binomials recover the bounded local numerical values.
The source majorant bounds its nonnegative local rule polynomial by
\(R^E<2^R<2^K\). Its divisibility by \(p^{Km}\ge2^K\)
therefore makes it zero. This produces every local row of the
numerical proof calculus.

In completeness use CRT at precision \(d_i^{2K}\), whose common
period is \(B^{2K}=N^2\). A falling-factorial numerator contains a
factor divisible by \(p^{2Km}\), while division by \(k!\) loses
at most \(k-1\le K-1\) powers of \(p\). At least \(Km\)
powers remain, so all range divisibilities hold. This argument needs
only the coprimality and product of the \(d_i\), both proved for the
descending moduli above. The semantic lift preserves the same precision.

Finally the local Horn-flow construction uses pairwise coprime weights
\(w(j)=1+jF\). Its beta-coded accumulator starts at \(w(1)=F+1\)
and ends at 1. Multiplying the local flow equations and taking prime
valuations recovers the multiplicity balance of conclusions and
premises. If contradiction label 1 were outside the closure of the
selected rules, every rule concluding outside the closure would have
an outside premise; the required excess of one outside conclusion is
impossible. Conversely a finite proof tree supplies the flow, with
padding to any sufficiently large horizon. This is the retained
numerical-calculus equivalence, followed by the two logical interfaces
stated in Theorem~\ref{thm:main}.

The order is important: frame and Pell indices; six exact powers;
first exact binomial; factorial; remaining binomial ranges; finite
local rows; refutation. Completeness chooses a finite proof first,
then its bounded semantic data, then the CRT representatives and their
lift, and only afterwards the much larger Pell witnesses.

\subsection{Literal compilation and independent review}

The finite graph exports to 339 natural coordinates and 288 quadratic
rows after affine normalization. A further 106 certified eliminations
leave the displayed 233 coordinates and 182 rows. Each inverse
substitution is an integer polynomial proved nonnegative on the natural
orthant; the two restorations with negative coefficients have exact
rational sum-of-squares certificates. Expansion gives exactly 7,313
monomials and degree four. All 233 variables occur.

An independent arithmetic review checks the unbounded floor, CRT, cap,
and horizon arguments. Its separate coefficient replay verifies the
entire expanded dictionary. Bounded controls exercise 2,197 factorial
cases, ten reciprocal products, and 2,452 cap tuples; they supplement
the proofs and are not used as evidence of universal unsolvability.
The files \code{natural/REVIEW.md}, \code{natural/INDEPENDENT_REPLAY.json},
and \code{natural/REPLAY.json} give the exact scopes.
